\section*{Bab \ref{ch:b2:reduction} --- Reduksi Endomorfisma}

\begin{solution}{exo:b2:reduction:1}
$A$: $\chi_A = X^2 - 2X - 3 = (X - 3)(X + 1)$. \omterm{def:b2:reduction:eigen}{Vektor eigennya}: untuk
$3$: $(1,1)$; untuk $-1$: $(1,-1)$. Jadi $P = \begin{pmatrix} 1 & 1\\ 1
& -1\end{pmatrix}$ memberi $P^{-1}AP = \operatorname{diag}(3, -1)$.

$B = J - I$ dengan $J$ matriks yang semua unsurnya satu. Matriks $J$
berank $1$ dengan $Jv = 3v$ untuk $v = (1,1,1)$ dan $Jw = 0$ pada bidang
$x + y + z = 0$: jadi \omterm{def:b2:reduction:eigen}{spektrum} $B$ adalah $\{2, -1\}$ dengan ruang eigen
$\operatorname{Vect}(1,1,1)$ (berdimensi $1$) dan $\{x + y + z = 0\}$
(berdimensi $2$, dengan basis $(1,-1,0), (1,0,-1)$). Maka $P$ yang
berkolomkan ketiganya memberi $P^{-1}BP = \operatorname{diag}(2, -1, -1)$.
\end{solution}

\begin{solution}{exo:b2:reduction:2}
\emph{Lewat ruang eigen:} $\chi_C = (X-1)^2$, dengan satu \omterm{def:b2:reduction:eigen}{nilai eigen} $1$;
$\ker(C - I) = \ker\begin{pmatrix} 0&1\\ 0&0\end{pmatrix}$ adalah
garis $\operatorname{Vect}(e_1)$, berdimensi $1 < 2 = m_1$, jadi tidak
\omterm{def:b2:reduction:diag}{dapat didiagonalkan} (\cref{thm:b2:reduction:diagcrit}).

\emph{Lewat \omterm{def:b2:reduction:polyu}{polinomial minimal}:} $\mu_C$ membagi $(X-1)^2$ dan $C \neq I$,
jadi $\mu_C = (X-1)^2$: akarnya ganda, sehingga tidak \omterm{def:b2:reduction:diag}{dapat didiagonalkan}
(\cref{cor:b2:reduction:minpolycrit}).
\end{solution}

\begin{solution}{exo:b2:reduction:3}
$X^2 - 5X + 6 = (X-2)(X-3)$: terurai lengkap dengan akar sederhana, jadi
$u$ \omterm{def:b2:reduction:diag}{dapat didiagonalkan} (\cref{cor:b2:reduction:minpolycrit}), dengan
$\operatorname{Sp}(u) \subseteq \{2, 3\}$. \omterm{def:b2:reduction:eigen}{Spektrum} yang mungkin:
$\{2\}$ (yaitu $u = 2\,\mathrm{id}$), $\{3\}$ (yaitu $u = 3\,\mathrm{id}$),
atau $\{2, 3\}$.

Pangkatnya: carilah $u^k = a_k\,\mathrm{id} + b_k\,u$. Pada \omterm{def:b2:reduction:eigen}{ruang eigennya},
syarat itu berbunyi $2^k = a_k + 2b_k$ dan $3^k = a_k + 3b_k$; setelah
diselesaikan, $b_k = 3^k - 2^k$ dan $a_k = 3\cdot2^k - 2\cdot 3^k$:
\[
u^k = (3\cdot 2^k - 2\cdot 3^k)\,\mathrm{id} + (3^k - 2^k)\, u .
\]
(Berlaku untuk ketiga \omterm{def:b2:reduction:eigen}{spektrumnya}: kesamaannya berlaku \omterm{def:b2:reduction:eigen}{nilai eigen} demi \omterm{def:b2:reduction:eigen}{nilai eigen}.)
\end{solution}

\begin{solution}{exo:b2:reduction:4}
Karena $u$ \omterm{def:b2:reduction:diag}{dapat didiagonalkan}, $P = \prod_{\lambda}(X - \lambda)$ atas
\omterm{def:b2:reduction:eigen}{spektrumnya} menganihilasi $u$, terurai lengkap, dan berakar sederhana.
Maka $P(u|_F) = P(u)|_F = 0$: pembatasannya dianihilasi oleh polinomial
terurai lengkap yang berakar sederhana, jadi ia \omterm{def:b2:reduction:diag}{dapat didiagonalkan}
(\cref{cor:b2:reduction:minpolycrit}).
\end{solution}

\begin{solution}{exo:b2:reduction:5}
Di sini $\chi_A = X^2 - X - 1$, dengan akar $\varphi$ dan $\psi$ yang
berbeda: jadi \omterm{def:b2:reduction:diag}{dapat didiagonalkan}, dengan \omterm{def:b2:reduction:eigen}{vektor eigen} $(\varphi, 1)$ dan
$(\psi, 1)$. Rekurensinya memberi $\begin{pmatrix} F_{n+1}\\ F_n
\end{pmatrix} = A^n \begin{pmatrix}1\\ 0\end{pmatrix}$. Uraikan $(1, 0)$
pada \omterm{def:b2:reduction:eigen}{vektor eigennya}: $(1,0) = \frac{1}{\varphi - \psi}\bigl((\varphi, 1)
- (\psi, 1)\bigr)$ dengan $\varphi - \psi = \sqrt5$. Menerapkan $A^n$
mengalikan tiap komponen eigennya dengan pangkat ke-$n$ \omterm{def:b2:reduction:eigen}{nilai eigennya};
setelah koordinat keduanya dibaca:
\[
F_n = \frac{\varphi^n - \psi^n}{\sqrt 5} .
\]
(Periksa: $n = 1$ memberi $\frac{\varphi - \psi}{\sqrt5} = 1$.)
\end{solution}

\begin{solution}{exo:b2:reduction:6}
Misalkan $P = \prod_i (X - \mu_i)$ menganihilasi $u^2$, terurai lengkap
dengan akar sederhana $\mu_i$ (yaitu \omterm{def:b2:reduction:eigen}{spektrum} $u^2$). Karena $u$ punya
invers, $0$ bukan \omterm{def:b2:reduction:eigen}{nilai eigen} $u^2$ (sebab $\det u^2 = (\det u)^2 \neq
0$), jadi semua $\mu_i \neq 0$. Maka
\[
Q(X) = \prod_i (X^2 - \mu_i) = \prod_i (X - \sqrt{\mu_i})(X +
\sqrt{\mu_i})
\]
menganihilasi $u$: $\;Q(u) = \prod_i (u^2 - \mu_i\,\mathrm{id}) =
P(u^2) = 0$. Akarnya $\pm
\sqrt{\mu_i}$ (yakni akar kuadrat kompleks) berbeda sepasang demi sepasang
karena $\mu_i$ berbeda dan tak nol (sebab $\sqrt{\mu_i} = -\sqrt{\mu_j}$
akan memberi $\mu_i = \mu_j$). Terurai lengkap dan berakar sederhana:
jadi $u$ \omterm{def:b2:reduction:diag}{dapat didiagonalkan}.

Contoh penyangkal tanpa keterbalikan: $u = \begin{pmatrix} 0 & 1\\
0 & 0\end{pmatrix}$: di sini $u^2 = 0$ \omterm{def:b2:reduction:diag}{dapat didiagonalkan}, $u$ tidak.
\end{solution}

\begin{solution}{exo:b2:reduction:7}
$A = 2I + N$ dengan $N$ geseran ($N e_2 = e_1$, $Ne_3 = e_2$),
$N^3 = 0$, $N^2 = E_{13}$: inilah \emph{tepat} \omterm{thm:b2:reduction:dunford}{penguraian Dunfordnya}
($2I$ diagonal, $N$ nilpoten, dan keduanya komutatif; ketunggalannya
menjadikannya satu-satunya). Binomial atas suku yang komutatif memberi
\[
A^k = 2^k I + k 2^{k-1} N + \binom k2 2^{k-2} N^2
= \begin{pmatrix}
2^k & k2^{k-1} & \binom k2 2^{k-2}\\
0 & 2^k & k2^{k-1}\\
0 & 0 & 2^k
\end{pmatrix},
\]
\[
\eu^{tA} = \eu^{2t}\Bigl(I + tN + \frac{t^2}{2}N^2\Bigr)
= \eu^{2t}\begin{pmatrix}
1 & t & t^2/2\\
0 & 1 & t\\
0 & 0 & 1
\end{pmatrix}.
\]
\end{solution}

\begin{solution}{exo:b2:reduction:8}
$X^k - 1$ menganihilasi $A$ dan terurai lengkap atas $\C$ dengan $k$ akar
berbeda $\eu^{2\iu\pi j/k}$: jadi $A$ \omterm{def:b2:reduction:diag}{dapat didiagonalkan}
(\cref{cor:b2:reduction:minpolycrit}) dan \omterm{def:b2:reduction:eigen}{nilai eigennya}, sebagai akar
$X^k - 1$, adalah akar ke-$k$ dari satu.

Jika lagi-lagi $A$ serupa dengan matriks segitiga berdiagonal
satu: semua \omterm{def:b2:reduction:eigen}{nilai eigennya} sama dengan $1$, dan $A$ yang \omterm{def:b2:reduction:diag}{dapat
didiagonalkan} dengan \omterm{def:b2:reduction:eigen}{nilai eigen} tunggal $1$ pastilah $P\,I\,P^{-1} = I$.
\end{solution}

\begin{solution}{exo:b2:reduction:9}
Tulis $E = \bigoplus_\lambda E_\lambda(u)$
(\cref{thm:b2:reduction:diagcrit}). Tiap $E_\lambda(u)$ stabil terhadap
$v$: untuk $x \in E_\lambda$ berlaku $u(v(x)) = v(u(x)) = \lambda
v(x)$. Pembatasan $v$ pada $E_\lambda(u)$ \omterm{def:b2:reduction:diag}{dapat didiagonalkan}
(\cref{exo:b2:reduction:4}): pilihlah basis $E_\lambda(u)$ yang terdiri
atas \omterm{def:b2:reduction:eigen}{vektor eigen} $v$. Menyambung basis-basis itu atas semua $\lambda$
memberi basis $E$ yang vektornya adalah \omterm{def:b2:reduction:eigen}{vektor eigen} bagi \emph{keduanya},
yakni $u$ (lewat keanggotaan di $E_\lambda(u)$) dan $v$ (menurut konstruksinya).
\end{solution}

\begin{solution}{exo:b2:reduction:10}
($\Rightarrow$) Misalkan $u$ \omterm{def:b2:reduction:diag}{dapat didiagonalkan} dan $F$ stabil. Maka
$u|_F$ \omterm{def:b2:reduction:diag}{dapat didiagonalkan} (\cref{exo:b2:reduction:4}): jadi $F$ punya
basis berisi \omterm{def:b2:reduction:eigen}{vektor eigen}, yang di dalam tiap ruang eigen global
$E_\lambda$ dapat diperluas menjadi basis $E_\lambda$ (lewat teorema
basis tak lengkap di dalam $E_\lambda$, berawal dari bagian basis $F$
yang terletak di sana --- perhatikan $F = \bigoplus_\lambda (F \cap
E_\lambda)$ sebab $u|_F$ \omterm{def:b2:reduction:diag}{dapat didiagonalkan}). Vektor yang ditambahkan
tadi merentang pelengkap yang stabil (masing-masing terletak di suatu
$E_\lambda$, jadi rentangnya stabil terhadap $u$).

($\Leftarrow$) Misalkan $F = \sum_\lambda E_\lambda(u)$ (sebuah subruang
stabil) dan $G$ pelengkapnya yang stabil. Jika $G \neq \{0\}$, maka
$\chi_{u|_G}$ terurai lengkap atas $\C$, sehingga $u|_G$ punya \omterm{def:b2:reduction:eigen}{vektor
eigen} $x \in G$ (\cref{thm:b2:reduction:trigonalization}, atau langsung
lewat keberadaan akarnya); tetapi setiap \omterm{def:b2:reduction:eigen}{vektor eigen} $u$ terletak di
$F$, jadi $x \in F \cap G = \{0\}$: kontradiksi. Karenanya $G = \{0\}$
dan $E = F$: \omterm{def:b2:reduction:eigen}{ruang eigennya} mengisi $E$, yakni $u$ \omterm{def:b2:reduction:diag}{dapat didiagonalkan}.
\end{solution}

\begin{solution}{exo:b2:reduction:11}
Trigonalkan: $A = PTP^{-1}$, $T = \begin{pmatrix} \lambda & c\\ 0 &
\mu\end{pmatrix}$, dengan $\abs\lambda, \abs\mu < 1$. Maka $A^k =
PT^kP^{-1}$, sehingga cukuplah menunjukkan $T^k \to 0$.

\emph{\omterm{def:b2:reduction:eigen}{Nilai eigen} berbeda:} induksi memberi
\[
T^k = \begin{pmatrix}
\lambda^k & c\,\dfrac{\lambda^k - \mu^k}{\lambda - \mu}\\[4pt]
0 & \mu^k
\end{pmatrix},
\]
dan tiap unsurnya menuju $0$ (sebab $\abs{\lambda}^k, \abs\mu^k \to 0$).

\emph{\omterm{def:b2:reduction:eigen}{Nilai eigen} sama ($\mu = \lambda$):} di sini $T = \lambda I + cE_{12}$
dan $T^k = \lambda^k I + k\lambda^{k-1}cE_{12}$; unsur
$k\lambda^{k-1} \to 0$ karena $\abs\lambda < 1$ (geometri mengalahkan
polinomial). Pada kedua kasusnya $T^k \to 0$ unsur demi unsur, sehingga
$A^k = PT^kP^{-1} \to 0$ (perkalian matriks oleh $P, P^{-1}$ yang tetap
bersifat kontinu pada unsurnya --- tiap unsur hasil kalinya adalah
kombinasi linear yang tetap).
\end{solution}

\begin{solution}{exo:b2:reduction:12}
Kernel $\ker u$ berdimensi $n - 1$ (rank--nulitas), jadi $0$ adalah
\omterm{def:b2:reduction:eigen}{nilai eigen} bermultiplisitas geometrik $n - 1$, dan $\chi_u$ habis
dibagi $X^{n-1}$
(\cref{def:b2:reduction:charpoly}: geometrik $\leq$ aljabar).
Tulis $\chi_u = X^{n-1}(X - \alpha)$; karena koefisien
$X^{n-1}$ adalah $-\operatorname{tr} u$, diperoleh $\alpha =
\operatorname{tr} u$, yakni $\chi_u = X^{n-1}(X - \operatorname{tr}
u)$.

Jika $\operatorname{tr} u \neq 0$: \omterm{def:b2:reduction:eigen}{nilai eigen}
$\operatorname{tr} u$ adalah akar $\chi_u$, jadi ia mengusung
\omterm{def:b2:reduction:eigen}{vektor eigen}; ruang eigen untuk $0$ dan untuk $\operatorname{tr} u$
berdimensi $n - 1$ dan $\geq 1$, yang berjumlah $\geq n$: keduanya
mengisi $E$, sehingga $u$ \omterm{def:b2:reduction:diag}{dapat didiagonalkan}
(\cref{thm:b2:reduction:diagcrit}). Jika $\operatorname{tr} u = 0$:
menurut \cref{exo:b2:linalg:5}, $u^2 = (\operatorname{tr} u)u = 0$
padahal $u \neq 0$, jadi $u$ nilpoten tak nol, sedangkan nilpoten yang
\omterm{def:b2:reduction:diag}{dapat didiagonalkan} pastilah nol (\cref{prop:b2:reduction:nilpotent}):
jadi tidak \omterm{def:b2:reduction:diag}{dapat didiagonalkan}.
\end{solution}

\begin{solution}{pb:b2:reduction:1}
\textbf{1.} Sebanyak $k - 1$ koordinat pertama pada $Cv_n$ adalah
$u_{n+1}, \dots, u_{n+k-1}$ (karena superdiagonalnya menggeser), dan yang
terakhir adalah $a_0 u_n + \dots + a_{k-1}u_{n+k-1}$. Jadi $v_{n+1} =
Cv_n$ berlaku untuk setiap $n$ bila dan hanya bila koordinat terakhirnya
cocok untuk setiap $n$, yakni bila dan hanya bila $(\mathcal R)$
berlaku. Dengan mengiterasikannya, $v_n = C^nv_0$.

\textbf{2.} Uraikan $D_k(X) = \det(XI_k - C)$ sepanjang kolom
pertama: kedua unsurnya yang tak nol adalah $X$ (pada kedudukan $(1,1)$)
dan $-a_0$ (pada kedudukan $(k,1)$). Minor pertamanya berbentuk
$D_{k-1}$ untuk koefisien $a_1, \dots, a_{k-1}$; minor keduanya
segitiga atas dengan diagonal $-1$, jadi \omterm{def:b2:linalg:det}{determinannya} $(-1)^{k-1}$,
dengan tanda $(-1)^{k+1}$ dari kedudukannya. Induksi pada $k$
(dengan pangkal $k = 1$: $X - a_0$) memberi
\[
D_k(X) = X\bigl(X^{k-1} - a_{k-1}X^{k-2} - \dots - a_1\bigr) -
a_0 = P(X).
\]
Untuk $\mu_C$: karena $Q(C^{\mathsf T}) = Q(C)^{\mathsf T}$ untuk setiap
polinomial, $C$ dan $C^{\mathsf T}$ punya penganihilasi yang sama, jadi
punya \omterm{def:b2:reduction:polyu}{polinomial minimal} yang sama. Untuk $C^{\mathsf T}$: kolomnya
berbunyi $C^{\mathsf T}e_1 = e_2$, \dots, $C^{\mathsf
T}e_{k-1} = e_k$, sehingga $(e_1, C^{\mathsf T}e_1, \dots, (C^{\mathsf
T})^{k-1}e_1)$ tak lain basis kanoniknya: bebas. Maka polinomial $Q
\neq 0$ berderajat $< k$ punya $Q(C^{\mathsf T})e_1 \neq 0$
(sebab ia kombinasi tak trivial atas vektor basis), jadi $\deg\mu \geq
k$. Karena $\mu \mid \chi = P$ dengan $\deg P = k$: $\mu_C = P$.

\textbf{3.} Untuk $v = (1, \lambda, \dots, \lambda^{k-1})^{\mathsf
T}$: baris $1$ sampai $k-1$ pada $Cv$ memberi $\lambda, \lambda^2, \dots,
\lambda^{k-1}$, yakni $\lambda$ kali $k - 1$ unsur pertama
$v$; sedangkan baris terakhirnya memberi $\sum_m a_m\lambda^m = \lambda^k -
P(\lambda) = \lambda^k = \lambda\cdot\lambda^{k-1}$. Jadi $Cv =
\lambda v$. Sebaliknya, persamaan $(Cx)_i = \lambda x_i$ untuk
$i < k$ berbunyi $x_{i+1} = \lambda x_i$: jadi setiap \omterm{def:b2:reduction:eigen}{vektor eigen}
sebanding dengan $v$ --- artinya tiap ruang eigen berdimensi tepat
$1$. \omterm{def:b2:reduction:diag}{Dapat didiagonalkan} bila dan hanya bila dimensi \omterm{def:b2:reduction:eigen}{ruang eigennya}
berjumlah $k$ (\cref{thm:b2:reduction:diagcrit}), bila dan hanya bila ada
$k$ \omterm{def:b2:reduction:eigen}{nilai eigen} berbeda, bila dan hanya bila $P$ punya $k$ akar berbeda
(sebab \omterm{def:b2:reduction:eigen}{nilai eigennya} adalah akar $\chi_C = P$).

\textbf{4.} Tiap $(\lambda_i^n)_n$ menyelesaikan $(\mathcal R)$, sebab
$\lambda_i^{n+k} = \lambda_i^n\,\lambda_i^k =
\lambda_i^n\sum_m a_m\lambda_i^m$. Kebebasannya: kombinasi yang nol,
$\sum_i c_i\lambda_i^n = 0$ untuk $n = 0, \dots, k-1$,
merupakan sistem Vandermonde (\cref{exo:b2:linalg:11}) pada $c_i$:
jadi semua $c_i = 0$. Ruang penyelesaiannya berdimensi $k$ (pertanyaan 6,
yang buktinya dasar dan berdiri sendiri): jadi $k$ penyelesaian bebas
membentuk basis, dan koordinatnya tunggal.

\textbf{5.} $P = X^2 - X - 6 = (X - 3)(X + 2)$: penyelesaian umumnya
$u_n = A\,3^n + B(-2)^n$. Syarat awalnya: $A + B = 1$, $3A -
2B = 8$, jadi $A = 2$ dan $B = -1$:
\[
u_n = 2\cdot 3^n - (-2)^n .
\]
(Periksa: $u_2 = u_1 + 6u_0 = 14$ dan $2\cdot9 - 4 = 14$.)

\textbf{6.} $P(S)\bigl((u_n)\bigr)$ adalah barisan $n \mapsto
u_{n+k} - a_{k-1}u_{n+k-1} - \dots - a_0u_n$: ia nol bila dan hanya bila
$(\mathcal R)$ berlaku, jadi himpunan penyelesaiannya adalah $\ker P(S)$,
sebuah subruang. Pemetaan $\ker P(S) \to \C^k$, $u \mapsto (u_0, \dots,
u_{k-1})$, bersifat linear, injektif (sebab rekurensinya menentukan
$u_{k}, u_{k+1}, \dots$ dari $k$ nilai pertamanya, secara induktif)
dan surjektif (tetapkan $u_n$ secara rekursif dari data awal apa pun):
jadi dimensinya $k$.

\textbf{7.} Bukti \cref{thm:b2:reduction:kernels} hanya memakai
kesamaan Bézout di $\C[X]$ serta kenyataan bahwa
polinomial dalam satu endomorfisma tetap saling komutatif. Tak satu pun
menyinggung dimensi ruang sekitarnya: jadi lemanya berlaku kata demi
kata bagi $S \in \mathcal{L}(\mathcal{S})$. Karenanya
\[
\ker P(S) = \bigoplus_{i=1}^{r}
\ker\,(S - \lambda_i\,\mathrm{id})^{m_i}.
\]

\textbf{8.} Untuk $Q \in \C[X]$: barisan $(S -
\lambda)\bigl(Q(n)\lambda^n\bigr)_n$ bersuku ke-$n$
$Q(n{+}1)\lambda^{n+1} - \lambda Q(n)\lambda^n =
\lambda^{n+1}(\Delta Q)(n)$, dengan $\Delta Q = Q(X{+}1) - Q(X)$
berderajat $\deg Q - 1$ (sebab suku utamanya saling coret). Dengan
mengiterasikannya, $(S - \lambda)^m\bigl(Q(n)\lambda^n\bigr) =
\bigl(\lambda^{n+m}(\Delta^m Q)(n)\bigr)_n$, dan $\Delta^m Q = 0$
bila $\deg Q \leq m - 1$: jadi himpunan di ruas kanan termuat di
kernelnya. Kernel itu subruang berdimensi $m$: barisan
$(n^j\lambda^n)_n$, $0 \leq j < m$, bersifat bebas, sebab
$\sum_j c_j n^j\lambda^n = 0$ untuk setiap $n$ memaksa (setelah dibagi
$\lambda^n \neq 0$) polinomial $\sum_j c_jX^j$ bernilai nol di
setiap $n \in \N$, jadi memaksanya menjadi nol. Sebaliknya $\dim\ker(S -
\lambda)^m \leq m$: setelah $(S - \lambda)^m = \sum_j
\binom mj(-\lambda)^{m-j}S^j$ diuraikan, persamaan $(S - \lambda)^m u =
0$ menjadi rekurensi linear berorde $m$ (dengan koefisien utama
$1$), sehingga $u$ ditentukan oleh $u_0, \dots, u_{m-1}$ seperti pada
pertanyaan 6. Kesamaan dimensinya merampungkan buktinya.

\textbf{9.} Gabungkan pertanyaan 7 dan 8: setiap penyelesaian terurai
secara tunggal sebagai jumlah unsur $\ker(S -
\lambda_i)^{m_i}$, yakni $u_n = \sum_i Q_i(n)\lambda_i^n$ dengan
$\deg Q_i \leq m_i - 1$; polinomial $Q_i$ itu tunggal karena
penguraiannya langsung dan, di dalam tiap sukunya,
koefisien $Q_i$ adalah koordinat pada basis
$(n^j\lambda_i^n)_j$ (pertanyaan 8). Pemeriksaan kewarasan pada dimensinya:
$\sum_i m_i = k$.

\textbf{10.} $P = X^2 - 4X + 4 = (X - 2)^2$: penyelesaiannya $(a +
bn)2^n$. Data awalnya: $a = 1$, $2(a + b) = 0$, jadi $b = -1$:
\[
u_n = (1 - n)\,2^n .
\]
Periksa: $u_2 = 4u_1 - 4u_0 = -4$, dan $(1 - 2)\cdot4 = -4$.

\textbf{11.} Tulis $u_n = \lambda_1^n\bigl(c_1 + \sum_{i\geq2}
c_i(\lambda_i/\lambda_1)^n\bigr)$; tiap rasionya bermodulus $< 1$,
jadi kurungnya menuju $c_1 \neq 0$: karenanya $u_n \sim c_1\lambda_1^n$.
Khususnya $u_n \neq 0$ untuk $n$ yang besar, dan
\[
\frac{u_{n+1}}{u_n} =
\lambda_1\,\frac{c_1 + o(1)}{c_1 + o(1)} \longrightarrow
\lambda_1 .
\]

\textbf{12.} Hitunglah:
\[
q(a_{n+1}, b_{n+1}) = (a_n + 2b_n)^2 - 2(a_n + b_n)^2
= -a_n^2 + 2b_n^2 = -q(a_n, b_n).
\]
Karena $q(a_0, b_0) = 1 - 2 = -1$, diperoleh $a_n^2 - 2b_n^2 = (-1)^{n+1}$.
Secara struktural: $q(a, b) = (a - \sqrt2\,b)(a + \sqrt2\,b)$ dan
pemetaan linear $M$ mengalikan faktor $a + \sqrt2 b$ dengan $1 +
\sqrt2$ serta faktor $a - \sqrt2 b$ dengan $1 - \sqrt2$ (hitunglah:
$a_{n+1} + \sqrt2 b_{n+1} = (1 + \sqrt2)(a_n + \sqrt2 b_n)$);
jadi hasil kalinya dikalikan $(1 + \sqrt2)(1 - \sqrt2) = -1 =
\det M$ pada tiap langkah.

\textbf{13.} Karena $a_n^2 - 2b_n^2 = (a_n - \sqrt2 b_n)(a_n +
\sqrt2 b_n) = (-1)^{n+1}$, berlaku
\[
\Bigl|\frac{a_n}{b_n} - \sqrt2\Bigr|
= \frac{\abs{a_n^2 - 2b_n^2}}{b_n(a_n + \sqrt2 b_n)}
= \frac{1}{b_n(a_n + \sqrt2 b_n)} \leq \frac1{2b_n^2},
\]
dengan memakai $a_n \geq b_n \geq 1$ (secara induktif keduanya naik)
sehingga $a_n + \sqrt2 b_n \geq (1 + \sqrt2)b_n \geq 2b_n$. \omterm{def:b2:reduction:eigen}{Nilai eigen}
$M$: $\chi_M = X^2 - 2X - 1$, dengan akar $1 \pm \sqrt2$; karena $(a_0,
b_0)$ punya komponen tak nol pada \omterm{def:b2:reduction:eigen}{vektor eigen} dominannya (semua
unsurnya positif), berlaku $b_n \sim c(1 + \sqrt2)^n$ dengan $c > 0$
(pertanyaan 11). Karenanya galatnya $\asymp (1 + \sqrt2)^{-2n} = (3 +
2\sqrt2)^{-n}$: peluruhan geometri dengan rasio $1/(3 + 2\sqrt2) = 3
- 2\sqrt2 \approx 0.172$.

\textbf{14.} (a) Dari pertanyaan 9: $\abs{u_n} \leq \sum_i
\abs{Q_i(n)}\abs{\lambda_i}^n \leq \bigl(\sum_i
\abs{Q_i(n)}\bigr)\rho^n$, dan tiap $\abs{Q_i(n)} \leq C_i
n^{m_i - 1} \leq C_i n^{m-1}$ untuk $n \geq 1$: jumlahkan konstantanya.
(b) Misalkan $\rho' = \max_{i \geq 2}\abs{\lambda_i} <
\abs{\lambda_1}$ dan $d = \deg Q_1$ dengan koefisien utama $c
\neq 0$. Maka $u_n = Q_1(n)\lambda_1^n + R_n$ dengan $\abs{R_n}
\leq Cn^{m-1}\rho'^n$, dan
\[
\frac{R_n}{Q_1(n)\lambda_1^n} = O\Bigl(n^{m-1-d}
\bigl(\rho'/\abs{\lambda_1}\bigr)^n\Bigr) \longrightarrow 0
\]
(sebab geometri mengalahkan polinomial). Jadi
\[
u_n \sim Q_1(n)\,\lambda_1^n \sim c\,n^d\lambda_1^n,
\qquad
\frac{u_{n+1}}{u_n} \longrightarrow \lambda_1
\quad\text{(sebab } Q_1(n{+}1)/Q_1(n) \to 1\text{)}.
\]
Periksa pada pertanyaan 10: untuk $u_n = (1-n)2^n$ rasionya adalah
\[
\frac{(-n)2^{n+1}}{(1-n)2^n} = 2\,\frac{-n}{1-n}
\longrightarrow 2 = \lambda_1 .
\]

\textbf{15.} Induksi pada $n$. Untuk $n = 1$, $A_{ij}$ mencacah
jalan berpanjang $1$. Langkahnya: jalan berpanjang $n + 1$ dari $i$ ke
$j$ adalah jalan berpanjang $n$ dari $i$ ke suatu simpul $\ell$
yang disusul rusuk $\ell j$:
\[
\#\{\text{jalan}\} = \sum_{\ell} (A^n)_{i\ell}A_{\ell j} =
(A^{n+1})_{ij}.
\]

\textbf{16.} Di sini $J = 3\Pi$ dengan $\Pi = J/3$ proyeksi pada
$\operatorname{Vect}(1,1,1)$ sepanjang bidang $x + y + z = 0$
(sebab $\Pi^2 = \Pi$ karena $J^2 = 3J$). Maka $A = J - I = 2\Pi - (I -
\Pi)$, dan karena $\Pi$ dan $I - \Pi$ merupakan proyeksi yang saling
melengkapi,
\[
A^n = 2^n\,\Pi + (-1)^n (I - \Pi),
\qquad\text{yakni}\qquad
(A^n)_{ij} = \frac{2^n}3 + (-1)^n\Bigl(\delta_{ij} -
\frac13\Bigr),
\]
sehingga diperoleh kedua rumus yang tertulis tadi. Pada $n = 2$: diagonalnya
$(4 + 2)/3 = 2$ (yaitu jalan $i \to \ell \to i$ lewat kedua tetangga
$\ell$); di luar diagonal $(4 - 1)/3 = 1$ (yaitu satu-satunya jalan $i \to
\ell \to j$ lewat simpul ketiga).

\textbf{17.} Misalkan $w_n^{(0)}, w_n^{(1)}$ mencacah kata yang sah dan
berpanjang $n$, berturut-turut yang berakhiran $0$ dan $1$. Menambahkan
satu huruf: angka $0$ boleh menyusul apa saja, angka $1$ hanya boleh
menyusul $0$:
\[
\begin{pmatrix} w_{n+1}^{(0)}\\ w_{n+1}^{(1)}\end{pmatrix}
= \begin{pmatrix} 1 & 1\\ 1 & 0\end{pmatrix}
\begin{pmatrix} w_n^{(0)}\\ w_n^{(1)}\end{pmatrix}.
\]
Setelah dijumlahkan, $w_{n+2} = w_{n+1} + w_n$ (atau: syaratkan huruf
pertamanya). Dengan $w_1 = 2$, $w_2 = 3$: secara induktif $w_n = F_{n+2}$
(sebab $F_3 = 2$, $F_4 = 3$, dengan rekurensi yang sama). Laju tumbuhnya:
akar $X^2 - X - 1$ adalah $\varphi > \abs\psi$
(\cref{exo:b2:reduction:5}), dan komponen $\varphi$-nya
tak nol (sebab $w_n$ positif dan $\psi^n \to 0$), jadi pertanyaan
11 memberi $w_{n+1}/w_n \to \varphi = \frac{1 + \sqrt5}2$.

\textbf{18.} Di sini $A = \begin{psmallmatrix} 0&1&0\\ 1&0&1\\
0&1&0\end{psmallmatrix}$. Periksa:
\[
A(1, \pm\sqrt2, 1)^{\mathsf T}
= (\pm\sqrt2, 2, \pm\sqrt2)^{\mathsf T}
= \pm\sqrt2\,(1, \pm\sqrt2, 1)^{\mathsf T},
\qquad
A(1, 0, -1)^{\mathsf T} = 0 :
\]
sehingga \omterm{def:b2:reduction:eigen}{nilai eigennya} $\sqrt2, -\sqrt2, 0$ ($= 2\cos\frac\pi4,
2\cos\frac{3\pi}4, 2\cos\frac\pi2$). Uraikan $e_1$ pada
basis eigennya lalu baca koordinat ketiganya, atau pakailah kesetangkupan:
dengan $v_\pm = (1, \pm\sqrt2, 1)$, $v_0 = (1, 0, -1)$, terperiksa $e_1
= \frac14 v_+ + \frac14 v_- + \frac12 v_0$, sehingga untuk $n \geq 1$
\[
(A^n)_{13} = \Bigl(\tfrac14(\sqrt2)^n v_+ +
\tfrac14(-\sqrt2)^n v_- + 0\Bigr)_{\!3}
= \frac{(\sqrt2)^n + (-\sqrt2)^n}{4},
\]
yang nol untuk $n$ ganjil (grafnya bipartit: kedua ujungnya berjarak
genap), dan $2\cdot 2^{n/2}/4 = 2^{n/2 - 1}$ untuk $n$ genap. Pada $n = 4$:
$2^{1} = 2$, yang cocok dengan kedua jalan $1\,2\,1\,2\,3$ dan
$1\,2\,3\,2\,3$.

\textbf{19.} Jalan tertutup berpanjang $n$ dari $i$ berjumlah
$(A^n)_{ii}$; menjumlahkannya atas $i$ memberi $\operatorname{tr}(A^n)$.
Setelah $A$ ditrigonalkan (atas $\C$), $A^n$ menjadi segitiga dengan
diagonal $\lambda_i^n$: jadi $\operatorname{tr}(A^n) = \sum_i\lambda_i^n$.
Untuk segitiganya: $\operatorname{tr}(A^n) = 3\,\frac{2^n + 2(-1)^n}3 =
2^n + 2(-1)^n = 2^n + (-1)^n + (-1)^n$, yakni \omterm{def:b2:reduction:eigen}{spektrum} $\{2, -1,
-1\}$, yang selaras dengan pertanyaan 16.

\textbf{20.} Kolom $W^{\mathsf T}$: $W^{\mathsf T}e_i =
e_{i-1}$ untuk $i \geq 1$ dan $W^{\mathsf T}e_0 = e_{k-1}$;
setelah dinamai ulang menurut urutan $e_0, e_1, \dots$ inilah persis
matriks pendamping bagi $X^k - 1$ (dengan $a_0 = 1$ dan $a_m = 0$ untuk yang lain).
Menurut pertanyaan 2: $\chi_{W} = \chi_{W^{\mathsf T}} = X^k - 1 =
\mu_{W}$. Akarnya $\omega^j$ ($j = 0, \dots, k-1$) adalah $k$
akar ke-$k$ dari satu yang berbeda: jadi $W$ \omterm{def:b2:reduction:diag}{dapat didiagonalkan}
(pertanyaan 3, atau \cref{exo:b2:reduction:8}: sebab $W^k = I$). \omterm{def:b2:reduction:eigen}{Vektor
eigennya}: $Wf_j = \sum_m \omega^{-jm}e_{m+1} =
\sum_{m'}\omega^{-j(m'-1)}e_{m'} = \omega^j f_j$.

\textbf{21.} Sirkulan adalah polinomial dalam $W$, dan polinomial
dalam satu matriks tetap saling komutatif. Tiap $f_j$ merupakan
\omterm{def:b2:reduction:eigen}{vektor eigen} bagi setiap pangkatnya: $W^m f_j = \omega^{jm}f_j$, jadi
\[
Cf_j = \sum_m c_m\omega^{jm} f_j = \widehat c(\omega^j)\,f_j :
\]
maka basis $(f_0, \dots, f_{k-1})$ (bebas: Vandermonde atas
$\omega^{-j}$ yang berbeda, \cref{exo:b2:linalg:11}) mendiagonalkan
setiap sirkulan sekaligus, dengan \omterm{def:b2:reduction:eigen}{nilai eigen} seperti dinyatakan tadi.

\textbf{22.} \omterm{def:b2:linalg:det}{Determinannya} adalah hasil kali \omterm{def:b2:reduction:eigen}{nilai eigennya}
(diagonalkan saja): $\det C = \prod_{j}\widehat c(\omega^j)$. Untuk $k
= 3$, dengan $c_0 = a$, $c_1 = b$, $c_2 = c$ dan $\omega = j =
\eu^{2\iu\pi/3}$:
\[
\det C = (a + b + c)(a + bj + cj^2)(a + bj^2 + cj^4),
\]
dan $j^4 = j$: jadi persis pemfaktoran pada
\cref{exo:b2:linalg:8}.

\textbf{23.} $M = \frac12(W + W^{-1})$ adalah sirkulan (sebab $W^{-1} =
W^{k-1}$), dengan \omterm{def:b2:reduction:eigen}{nilai eigen} $\frac12(\omega^j + \omega^{-j}) =
\cos\frac{2\pi j}k$ pada basis $f_j$ yang sama. Koordinatnya: tulis
$x^{(0)} = \sum_j \alpha_j f_j$. Koordinat $f_j$ berjumlah
$\sum_m \omega^{-jm}$, yang bernilai $k$ untuk $j = 0$ dan $0$
selain itu (jumlah geometri dengan rasio $\omega^{-j} \neq 1$).
Menjumlahkan koordinat $x^{(0)}$: $\sum_m x^{(0)}_m =
\alpha_0\,k$, jadi $\alpha_0 = \frac1k\sum_m x^{(0)}_m$, yakni rata-ratanya.

\textbf{24.} Di sini $x^{(n)} = M^nx^{(0)} = \sum_j
\alpha_j\cos^n\bigl(\tfrac{2\pi j}k\bigr)f_j$. Untuk $k$ ganjil,
$\abs{\cos(2\pi j/k)} < 1$ bagi setiap $j \neq 0$ (sebab sudutnya
tak pernah $0$ atau $\pi$), jadi semua sukunya kecuali $j = 0$ menuju $0$:
$x^{(n)} \to \alpha_0 f_0$, yakni vektor konstan yang sama dengan
rata-ratanya --- perataan pada cincin ganjil menyeragamkan. Untuk $k = 4$
\omterm{def:b2:reduction:eigen}{nilai eigennya} $1, 0, -1, 0$: suku $j = 2$, yakni
$\alpha_2(-1)^nf_2$ dengan $f_2 = (1, -1, 1, -1)^{\mathsf T}$,
berayun selamanya. Halangannya adalah rata-rata berselang-seling:
setelah koordinat $x^{(0)}$ dikalikan $(-1)^m$ lalu
dijumlahkan, perhitungan jumlah geometri yang sama memberi $\sum_m
(-1)^mx^{(0)}_m = 4\alpha_2$: jadi prosesnya konvergen bila dan hanya bila
$x^{(0)}_0 - x^{(0)}_1 + x^{(0)}_2 - x^{(0)}_3 = 0$, dan ketika itu
ia konvergen ke rata-ratanya.

\textbf{25.} Matriks pendamping mengubah rekurensi skalar
berorde $k$ menjadi rekurensi vektor berorde satu, sehingga
rumus tertutupnya menjadi pernyataan tentang $C^n$ --- yakni
kandang reduksi sendiri (pertanyaan 1--5). Lema penguraian kernel
murni aljabar polinomial (Bézout ditambah kekomutatifan), jadi ia
memecah $\ker P(S)$ walaupun $\mathcal{S}$ berdimensi tak hingga
(pertanyaan 7--9). \omterm{def:b2:reduction:eigen}{Nilai eigen} dominan menguasai laju tumbuh karena
setiap sumbangan lain terabaikan secara geometri setelah
dinormalkan --- dan itu pula sebabnya galat Pell meluruh secepat
kuadrat akar dominannya (pertanyaan 11--14). Pangkat
matriks ketetanggaan mencacah jalan karena perkalian matriks menjumlah
atas simpul antaranya, sehingga \omterm{def:b2:reduction:eigen}{spektrumnya} mencacah jalan tertutup
(pertanyaan 15--19). Matriks yang komutatif berbagi basis eigen, dan
satu basis Fourier lalu mendiagonalkan seluruh aljabar sirkulan
dalam satu tarikan (pertanyaan 20--24). Puncaknya: teorema dasar
rekurensi linear (pertanyaan 9); dan untuk matriks tak negatif,
alasan mengapa akar dominan seperti $\varphi$ atau $1 +
\sqrt2$ otomatis real, positif dan sederhana adalah
teorema Perron--Frobenius, yang dibuktikan pada jilid Tahun ke-3.

\end{solution}
